\documentclass[aspectratio=169, xcolor={table,xcdraw}]{beamer}

\usepackage[T1]{fontenc}
\usepackage[utf8]{inputenc}
\usepackage{tgpagella}
\usepackage{tgcursor}
\usefonttheme{serif}
\usepackage[spanish, es-noshorthands]{babel}
\renewcommand{\tablename}{Tabla}

\usepackage{amsmath, amssymb, bm}
\usepackage{graphicx}
\usepackage{booktabs}
\usepackage[most]{tcolorbox}
\usepackage{tikz}
\usetikzlibrary{shapes.geometric, arrows.meta, positioning, calc, angles, quotes}

\usetheme{Madrid}
\usecolortheme{default}
\definecolor{ucbblue}{RGB}{0, 51, 102}
\setbeamercolor{structure}{fg=ucbblue}
\setbeamercolor{title}{fg=white, bg=ucbblue}
\setbeamercolor{frametitle}{fg=white, bg=ucbblue}

\title[Consolidación y Singularidades]{Consolidación del Jacobiano y Transición a Singularidades}
\subtitle{De la Velocidad Local a la Pérdida de Rango}
\author[B. Quiroga Turdera]{M.Sc. Ing. Bernardo Quiroga Turdera}
\institute[UCB Tarija]{Universidad Católica Boliviana ``San Pablo'' — Sede Tarija}
\date{Semana 8 - Sesión 2}

\begin{document}

\begin{frame} 
    \titlepage
\end{frame}

\begin{frame}{Problema Diagnóstico Inicial: Construcción Independiente[cite: 12]}
    \begin{columns}[T]
        \begin{column}{0.48\textwidth}
            \small
            Considere el manipulador planar 2R con[cite: 12]:
            $a_1 = a_2 = 1, \quad q_1 = 0, \quad q_2 = \frac{\pi}{2}$
            
            \vspace{0.2cm}
            Los orígenes y ejes en la base $\{0\}$ son[cite: 12]:
            \begin{align*}
                {}^0\mathbf{o}_0 &= [0, 0, 0]^T \\
                {}^0\mathbf{o}_1 &= [1, 0, 0]^T \\
                {}^0\mathbf{o}_2 &= [1, 1, 0]^T \\
                {}^0\mathbf{z}_0 &= {}^0\mathbf{z}_1 = [0, 0, 1]^T
            \end{align*}
            \textbf{Misión:} Construya cada columna, ensamble la matriz Jacobiana $J \in \mathbb{R}^{6 \times 2}$ e interprete cada columna físicamente[cite: 12].
        \end{column}
        \begin{column}{0.52\textwidth}
            \centering
            \begin{tikzpicture}[scale=1.4, thick]
                % Ejes cartesianos
                \draw[->, gray] (-0.5,0) -- (2,0) node[right]{$x_0$};
                \draw[->, gray] (0,-0.5) -- (0,2) node[above]{$y_0$};
                
                \coordinate (O0) at (0,0);
                \coordinate (O1) at (1,0);
                \coordinate (O2) at (1,1);
                
                % Eslabones
                \draw[gray!60, line width=3pt] (O0) -- (O1) node[midway, below=3pt, text=black]{$a_1$};
                \draw[gray!60, line width=3pt] (O1) -- (O2) node[midway, right=3pt, text=black]{$a_2$};
                
                % Vectores de Radio r_i = o_n - o_{i-1}
                \draw[-Stealth, green!60!black, line width=1.5pt] (O0) -- (O2) node[midway, above left]{$\mathbf{r}_1 = \mathbf{o}_2 - \mathbf{o}_0$};
                
                % Transformamos canvas para no superponer r2 sobre el eslabon a2
                \draw[-Stealth, orange, line width=1.5pt, transform canvas={xshift=-3pt}] (O1) -- (O2) node[midway, left]{\small $\mathbf{r}_2 = \mathbf{o}_2 - \mathbf{o}_1$};
                
                % Puntos de origen
                \filldraw[black] (O0) circle (1.5pt) node[below left]{$\mathbf{o}_0$};
                \filldraw[black] (O1) circle (1.5pt) node[below right]{$\mathbf{o}_1$};
                \filldraw[red] (O2) circle (1.5pt) node[above right]{$\mathbf{o}_2$ (Efector)};
            \end{tikzpicture}
        \end{column}
    \end{columns}
\end{frame}

\begin{frame}{Consolidación: Anatomía de la Matriz Jacobiana Geométrica[cite: 12]}
    El Jacobiano mapea velocidades articulares a velocidad espacial del efector[cite: 12]:
    \begin{center}
        $\mathbf{v}_e = \begin{bmatrix} \mathbf{v}_L \\ \boldsymbol{\omega} \end{bmatrix} = J(\mathbf{q})\dot{\mathbf{q}}$
    \end{center}

    \begin{columns}[T]
        \begin{column}{0.48\textwidth}
            \begin{tcolorbox}[colback=ucbblue!5, colframe=ucbblue, title=Junta Revoluta $i$, left=1mm, right=1mm, top=1mm, bottom=1mm]
                \small
                Genera velocidad angular, y velocidad tangencial debido al brazo de palanca[cite: 12]:
                \begin{equation*}
                    J_i = \begin{bmatrix} \mathbf{z}_{i-1} \times (\mathbf{o}_n - \mathbf{o}_{i-1}) \\ \mathbf{z}_{i-1} \end{bmatrix}
                \end{equation*}
            \end{tcolorbox}
        \end{column}
        \begin{column}{0.48\textwidth}
            \begin{tcolorbox}[colback=red!5, colframe=red!70!black, title=Junta Prismática $i$, left=1mm, right=1mm, top=1mm, bottom=1mm]
                \small
                Genera traslación pura a lo largo de su eje, sin inducir velocidad angular[cite: 12]:
                \begin{equation*}
                    J_i = \begin{bmatrix} \mathbf{z}_{i-1} \\ \mathbf{0} \end{bmatrix}
                \end{equation*}
            \end{tcolorbox}
        \end{column}
    \end{columns}
    
    \vfill
    \begin{alertblock}{Regla Estricta de Marcos[cite: 12]}
        \small Todos los vectores ($\mathbf{z}$ y $\mathbf{o}$) en una columna deben estar expresados en el \textbf{mismo marco de referencia} (usualmente la base $\{0\}$) antes del producto cruz.
    \end{alertblock}
\end{frame}

\begin{frame}{Resolución del Problema Diagnóstico: Junta 1 (1/3)[cite: 12]}
    Calculamos los componentes de velocidad lineal ($\mathbf{v}_{Li} = \mathbf{z}_{i-1} \times \mathbf{r}_i$) y angular ($\boldsymbol{\omega}_i = \mathbf{z}_{i-1}$)[cite: 12].

    \vfill
    \textbf{Para la Junta 1 (Revoluta):}
    \begin{itemize}
        \item \textbf{Vector Radio:} $\mathbf{r}_1 = \mathbf{o}_2 - \mathbf{o}_0 = [1, 1, 0]^T$
        \item \textbf{Velocidad Lineal:} $\mathbf{v}_{L1} = \mathbf{z}_0 \times \mathbf{r}_1 = [0, 0, 1]^T \times [1, 1, 0]^T = [-1, 1, 0]^T$
        \item \textbf{Columna 1:} $J_1 = \begin{bmatrix} \mathbf{v}_{L1}^T & \mathbf{z}_0^T \end{bmatrix}^T = [-1, 1, 0, 0, 0, 1]^T$
    \end{itemize}
    \vfill
\end{frame}

\begin{frame}{Resolución del Problema Diagnóstico: Junta 2 (2/3)[cite: 12]}
    \textbf{Para la Junta 2 (Revoluta):}

    \vfill
    \begin{itemize}
        \item \textbf{Vector Radio:} $\mathbf{r}_2 = \mathbf{o}_2 - \mathbf{o}_1 = [1, 1, 0]^T - [1, 0, 0]^T = [0, 1, 0]^T$
        \item \textbf{Velocidad Lineal:} $\mathbf{v}_{L2} = \mathbf{z}_1 \times \mathbf{r}_2 = [0, 0, 1]^T \times [0, 1, 0]^T = [-1, 0, 0]^T$
        \item \textbf{Columna 2:} $J_2 = \begin{bmatrix} \mathbf{v}_{L2}^T & \mathbf{z}_1^T \end{bmatrix}^T = [-1, 0, 0, 0, 0, 1]^T$
    \end{itemize}
    \vfill
\end{frame}

\begin{frame}{Resolución del Problema Diagnóstico: Matriz Final (3/3)[cite: 12]}
    Ensamblando la matriz final concatenando las columnas calculadas[cite: 12]:
    \begin{equation*}
        \boxed{ J = \begin{bmatrix} | & | \\ J_1 & J_2 \\ | & | \end{bmatrix} = \begin{bmatrix} -1 & -1 \\ 1 & 0 \\ 0 & 0 \\ 0 & 0 \\ 0 & 0 \\ 1 & 1 \end{bmatrix} } 
    \end{equation*}

    \vspace{0.4cm}
    Si solo nos interesa el movimiento de traslación en el plano XY, extraemos el \textbf{sub-Jacobiano planar} $J_p$ (las dos primeras filas)[cite: 12]:
    \begin{equation*}
        \boxed{ J_p = \begin{bmatrix} -1 & -1 \\ 1 & 0 \end{bmatrix} }
    \end{equation*}
\end{frame}

\begin{frame}{Mapeo de Velocidades y Significado de las Columnas[cite: 12]}
    El Jacobiano combina las contribuciones de velocidad de cada junta mediante una suma ponderada (combinación lineal)[cite: 12]. Dadas las velocidades articulares $\dot{\mathbf{q}} = [1, -1]^T \mathrm{rad/s}$[cite: 12]:
    
    {\footnotesize
    \begin{equation*}
        \mathbf{v}_e = J\dot{\mathbf{q}} = J_1 \dot{q}_1 + J_2 \dot{q}_2 = \begin{bmatrix} -1 \\ 1 \\ 0 \\ 0 \\ 0 \\ 1 \end{bmatrix}(1) + \begin{bmatrix} -1 \\ 0 \\ 0 \\ 0 \\ 0 \\ 1 \end{bmatrix}(-1) = \boxed{ \begin{bmatrix} 0 \\ 1 \\ 0 \\ 0 \\ 0 \\ 0 \end{bmatrix} } \text{[cite: 12]}
    \end{equation*}
    }
    
    \textbf{Interpretación Física[cite: 12]:}
    \begin{itemize}
        \item Cada columna representa la velocidad que se induciría en el efector si esa junta rotara a $1\,\mathrm{rad/s}$ y las demás estuvieran fijas[cite: 12].
        \item En este mapeo específico, las contribuciones angulares de ambas juntas se cancelan ($1 - 1 = 0$), dejando al efector en \textbf{traslación instantánea pura} a lo largo de $+y$[cite: 12].
    \end{itemize}
\end{frame}

\begin{frame}{Cinemática Inversa: Diferencial vs. Analítica[cite: 12]}
    \begin{center}
        \renewcommand{\arraystretch}{1.2}
        \footnotesize
        \begin{tabular}{|p{0.45\textwidth}|p{0.45\textwidth}|}
        \hline
        \rowcolor{gray!20}\textbf{IK Analítica / Global} & \textbf{IK Diferencial / Local} \\ \hline
        Pose deseada $\to$ Posiciones articulares ($T_d \to \mathbf{q}$) & Velocidad cartesiana $\to$ Vel. articulares ($\dot{\mathbf{x}}_d \to \dot{\mathbf{q}}$) \\ \hline
        Problema a nivel de pose global. & Relación local en la configuración actual. \\ \hline
        Produce múltiples soluciones discretas (Ramas). & Mapeo lineal dependiente de la configuración. \\ \hline
        Razonamiento geométrico / Algoritmo de Pieper. & Inversión del Jacobiano (cuando es admisible). \\ \hline
        \end{tabular}
    \end{center}
    
    \vfill
    Para un Jacobiano $J$ \textbf{cuadrado y no singular}, la IK diferencial es[cite: 12]:
    \begin{center}
        $\boxed{\dot{\mathbf{q}} = J^{-1}(\mathbf{q})\dot{\mathbf{x}}_d}$
    \end{center}
    
    \vfill
    \begin{alertblock}{No reemplaza a Pieper[cite: 12]}
        \footnotesize $J^{-1}$ resuelve cómo mantener una trayectoria local continua, pero \textbf{no} resuelve el problema global de alcanzar una pose desde cero.
    \end{alertblock}
\end{frame}

\begin{frame}{IK Diferencial: Ejemplo con el Brazo 2R[cite: 12]}
    Utilizando el sub-Jacobiano de traslación del problema diagnóstico[cite: 12]:
    \begin{equation*}
        J_p = \begin{bmatrix} -1 & -1 \\ 1 & 0 \end{bmatrix}
    \end{equation*}
    
    Invertimos la matriz $2 \times 2$[cite: 12]:
    \begin{equation*}
        J_p^{-1} = \frac{1}{\det(J_p)} \begin{bmatrix} 0 & 1 \\ -1 & -1 \end{bmatrix} = \frac{1}{(0 - (-1))} \begin{bmatrix} 0 & 1 \\ -1 & -1 \end{bmatrix} = \begin{bmatrix} 0 & 1 \\ -1 & -1 \end{bmatrix}
    \end{equation*}
    
    Si deseamos una velocidad instantánea en el efector final de $\dot{\mathbf{p}}_d = [0, 1]^T$[cite: 12]:
    \begin{equation*}
        \boxed{ \dot{\mathbf{q}} = J_p^{-1} \dot{\mathbf{p}}_d = \begin{bmatrix} 0 & 1 \\ -1 & -1 \end{bmatrix} \begin{bmatrix} 0 \\ 1 \end{bmatrix} = \begin{bmatrix} 1 \\ -1 \end{bmatrix} \mathrm{rad/s} } \text{[cite: 12]}
    \end{equation*}
\end{frame}

\begin{frame}{Rango del Jacobiano y Concepto de Singularidad[cite: 12]}
    \small
    \begin{columns}[T]
        \begin{column}{0.48\textwidth}
            El \textbf{rango} de una matriz es el número de columnas linealmente independientes. Físicamente, mide la cantidad de \textbf{direcciones instantáneas independientes} en las que el robot puede moverse[cite: 12].
            
            \vspace{0.2cm}
            Para $J \in \mathbb{R}^{m \times n}$, el rango máximo es:
            $\operatorname{rank}_{\max} = \min(m,n)$[cite: 12].
            
            \vspace{0.2cm}
            Una configuración es \textbf{singular} si[cite: 12]:
            \begin{equation*}
                \boxed{\operatorname{rank}(J(\mathbf{q})) < \operatorname{rank}_{\max}}
            \end{equation*}
        \end{column}
        \begin{column}{0.48\textwidth}
            \textbf{Ejemplo: Brazo 2R Estirado ($q_1=0, q_2=0$)}[cite: 12]
            \begin{equation*}
                J_p = \begin{bmatrix} 0 & 0 \\ 2 & 1 \end{bmatrix} \text{[cite: 12]}
            \end{equation*}
            Las dos columnas son proporcionales. 
            \begin{equation*}
                \boxed{\operatorname{rank}(J_p) = 1} \text{[cite: 12]}
            \end{equation*}
            \textbf{Interpretación:} El efector ha perdido la capacidad de generar velocidad instantánea independiente en la dirección $X$ (radial)[cite: 12].
        \end{column}
    \end{columns}
\end{frame}

\begin{frame}{El Criterio del Determinante (Solo Matrices Cuadradas)[cite: 12]}
    \small
    Para un Jacobiano cuadrado, $\det(J) = 0$ es una prueba de singularidad, pero \textbf{la definición universal es la pérdida de rango}[cite: 12].
    
    Derivación de $\det(J_p)$ para el 2R general[cite: 12]:
    \begin{equation*}
        J_p = \begin{bmatrix} -a_1 s_1 - a_2 s_{12} & -a_2 s_{12} \\ a_1 c_1 + a_2 c_{12} & a_2 c_{12} \end{bmatrix}
    \end{equation*}
    Calculamos el determinante y factorizamos términos comunes[cite: 12]:
    \begin{align*}
        \det(J_p) &= (-a_1 s_1 - a_2 s_{12})(a_2 c_{12}) - (-a_2 s_{12})(a_1 c_1 + a_2 c_{12}) \\
        &= a_1 a_2 (s_{12} c_1 - c_{12} s_1) \text{[cite: 12]}
    \end{align*}
    Usando la identidad trigonométrica $\sin(A-B) = \sin A \cos B - \cos A \sin B$[cite: 12]:
    \begin{equation*}
        \boxed{\det(J_p) = a_1 a_2 \sin q_2} \text{[cite: 12]}
    \end{equation*}
    La singularidad ocurre cuando $q_2 = 0$ o $q_2 = \pi$[cite: 12].
\end{frame}

\begin{frame}{Geometría de la Singularidad: Regular vs. Singular[cite: 12]}
    \begin{columns}[T]
        \begin{column}{0.5\textwidth}
            \textbf{Regular ($q_2 = \pi/2$)}[cite: 12]
            
            \centering
            \begin{tikzpicture}[scale=1.0, thick]
                \draw[->, gray] (0,0) -- (1.5,0);
                \draw[->, gray] (0,0) -- (0,1.5);
                
                \coordinate (O0) at (0,0);
                \coordinate (O1) at (1,0);
                \coordinate (O2) at (1,1);
                
                \draw[ucbblue, line width=2pt] (O0) -- (O1);
                \draw[ucbblue!70, line width=2pt] (O1) -- (O2);
                \filldraw[black] (O0) circle (1pt); \filldraw[black] (O1) circle (1pt); \filldraw[red] (O2) circle (1pt);
                
                % Velocidades inducidas
                \draw[-Stealth, red!70!black, line width=1.5pt] (O2) -- ++(-0.5, 0.5) node[above]{$J_{p1}$};
                \draw[-Stealth, orange, line width=1.5pt] (O2) -- ++(-0.5, 0) node[left]{$J_{p2}$};
            \end{tikzpicture}
            
            Columnas independientes. $\operatorname{rank}(J_p)=2$[cite: 12].
        \end{column}
        \begin{column}{0.5\textwidth}
            \textbf{Singular ($q_2 = 0$)}[cite: 12]
            
            \centering
            \begin{tikzpicture}[scale=1.0, thick]
                \draw[->, gray] (0,0) -- (2.5,0);
                \draw[->, gray] (0,0) -- (0,1);
                
                \coordinate (O0) at (0,0);
                \coordinate (O1) at (1,0);
                \coordinate (O2) at (2,0);
                
                \draw[ucbblue, line width=2pt] (O0) -- (O1);
                \draw[ucbblue!70, line width=2pt] (O1) -- (O2);
                \filldraw[black] (O0) circle (1pt); \filldraw[black] (O1) circle (1pt); \filldraw[red] (O2) circle (1pt);
                
                % Velocidades inducidas
                \draw[-Stealth, red!70!black, line width=1.5pt] (O2) -- ++(0, 1) node[right]{$J_{p1}$};
                \draw[-Stealth, orange, line width=1.5pt] (O2) -- ++(0, 0.5) node[left]{$J_{p2}$};
            \end{tikzpicture}
            
            Columnas paralelas. $\operatorname{rank}(J_p)=1$[cite: 12]. \\ Movimiento radial en X perdido[cite: 12].
        \end{column}
    \end{columns}
\end{frame}

\begin{frame}{Comportamiento de Velocidad Cerca de la Singularidad[cite: 12]}
    El problema de la singularidad no solo ocurre cuando el determinante es cero. \textbf{Cerca de la singularidad}, las velocidades articulares se amplifican dramáticamente[cite: 12].
    
    Ilustración con un Jacobiano diagonal simple dependiente de $\varepsilon > 0$[cite: 12]:
    \begin{equation*}
        J_\varepsilon = \begin{bmatrix} 1 & 0 \\ 0 & \varepsilon \end{bmatrix} \quad \implies \quad J_\varepsilon^{-1} = \begin{bmatrix} 1 & 0 \\ 0 & 1/\varepsilon \end{bmatrix} \text{[cite: 12]}
    \end{equation*}
    Para una velocidad cartesiana deseada finita $\dot{\mathbf{x}}_d = [0, 1]^T$, la velocidad articular requerida es[cite: 12]:
    \begin{equation*}
        \dot{\mathbf{q}} = J_\varepsilon^{-1} \dot{\mathbf{x}}_d = \begin{bmatrix} 0 \\ 1/\varepsilon \end{bmatrix} \text{[cite: 12]}
    \end{equation*}

    \begin{center}
        \begin{tabular}{ccc}
        \toprule
        $\varepsilon$ (Proximidad singular) & $\dot{\mathbf{x}}_d$ & $\dot{q}_2$ Requerida \\ \midrule
        $0.5$ & $1.0$ & $2$ \\
        $0.1$ & $1.0$ & $10$ \\
        $0.01$ & $1.0$ & $100$ \\ \bottomrule
        \end{tabular}
    \end{center}
\end{frame}

\begin{frame}{Exit Ticket[cite: 12]}
    \begin{tcolorbox}[colback=white, colframe=ucbblue, title=Verificación de Consolidación]
        \begin{enumerate}
            \item ¿Qué representa físicamente una sola columna de la matriz Jacobiana geométrica?[cite: 12]
            \item ¿Por qué todos los vectores en la fórmula $J_i = \mathbf{z}_{i-1} \times (\mathbf{o}_n - \mathbf{o}_{i-1})$ deben estar expresados en el mismo marco de referencia?[cite: 12]
            \item ¿En qué se diferencia la cinemática inversa diferencial de la analítica (global)?[cite: 12]
            \item Físicamente, ¿qué significa que un robot pierda rango en su Jacobiano?[cite: 12]
            \item ¿Por qué $\det(J) = 0$ no es la definición universal de una singularidad?[cite: 12]
            \item ¿Por qué las velocidades articulares se disparan al acercarse a una singularidad?[cite: 12]
        \end{enumerate}
    \end{tcolorbox}
\end{frame}

\end{document}
